\begin{table}[tbp]
    \centering
    \small
    \caption{Cinética fotoquímica e fração de poros abertos $O(t)$ no Experimento 2.}
    \label{tab:resultados_exp2_cinetica}
    \begin{tabular}{llccc}
        \toprule
        \textbf{Opsina} & \textbf{Padrão}    & \textbf{$\tau_{\text{off}}$}                 & \textbf{$O_{\text{médio}}$} & \textbf{$O_{\text{mín}}$} \\
        \midrule
        \multirow{2}{*}{ChR2}
                        & Fixa ($MS_f$)      & \multirow{2}{*}{\qty{13{,}0}{\milli\second}} & $0{,}572$                   & $0{,}180$                 \\
                        & Aleatória ($MS_r$) &                                              & $0{,}548$                   & $0{,}021$                 \\
        \midrule
        \multirow{2}{*}{Chronos}
                        & Fixa ($MS_f$)      & \multirow{2}{*}{\qty{3{,}6}{\milli\second}}  & $0{,}285$                   & $0{,}002$                 \\
                        & Aleatória ($MS_r$) &                                              & $0{,}279$                   & $0{,}000$                 \\
        \bottomrule
    \end{tabular}
    \vspace{0.2em}
    \begin{minipage}{0.95\textwidth}
        \footnotesize \textbf{Nota:} Valores na janela de estimulação ativa ($\qty{140}{\milli\second}$ por ciclo). A ChR2 sustenta patamar residual elevado ($O_{\text{mín}} = 0{,}180$), enquanto a Chronos atinge fechamento completo entre pulsos adjacentes ($O_{\text{mín}} \approx 0$).
    \end{minipage}
\end{table}
